\chapter{Greedy Algorithms}

\section{Introduction}

\term{Greedy algorithms} are a class of algorithms that make locally optimal choices at each step in order to find a global optimum. They are often used to solve optimization problems.

\begin{listu}
    \item \textbf{Goal:} find a solution $x$ maximizing or minimizing some objective function $f$. 
    \item \textbf{Challenge:} space of possible solutions $x$ is too large to search exhaustively.
    \item \textit{Insight:} $x$ is composed of several parts (e.g., $x$ is a set or a sequence).
    \item \textbf{Approach:} instead of computing $x$ directly, compute $x$ one part at a time.
    \begin{listu}
        \item Select the next part ``greedily'' to get the most immediate ``benefit'', which needs to be defined carefully for each problem.
        \item Polynomial running time is typically guaranteed.
        \item Need to prove that this will always return an optimal solution despite having no global view of the problem.
    \end{listu}
\end{listu}

\section{Interval Scheduling}

\subsection{Problem Definition}

\begin{listu}
    \item \textbf{Problem}
    
    \begin{listu}
        \item Job $j$ starts at time $s_j$ and finishes at time $f_j$.
        \item Two jobs $i$ and $j$ are compatible if $[s_i, f_i)$ and $[s_j, f_j)$ do not overlap.
        \item \textbf{Goal:} find a maximum-size subset of mutually compatible jobs.
    \end{listu}

    \item \textbf{Applications}

    \begin{listu}
        \item Scheduling jobs on a single machine.
        \item Scheduling classes in a classroom.
        \item Scheduling packets on a link.
    \end{listu}
\end{listu}

% TODO: counter examples

\begin{minipage}[t]{0.55\linewidth}
    \begin{center} \begin{tikzpicture}[baseline=(current bounding box.center)]
        \definecolor{lightGray}{gray}{0.9}
        \definecolor{darkGray}{gray}{0.7}

        \node[draw=none, fill=lightGray, minimum width=1cm, minimum height=0.5cm] at (0, 0) {};
        \node[draw=none, fill=lightGray, minimum width=1cm, minimum height=0.5cm] at (1.5, 0) {};
        \node[draw=none, fill=lightGray, minimum width=1cm, minimum height=0.5cm] at (3, 0) {};
        \node[draw=none, fill=lightGray, minimum width=1cm, minimum height=0.5cm] at (4.5, 0) {};

        \node[draw=none, fill=darkGray, minimum width=6cm, minimum height=0.5cm] at (2.25, -0.75) {};
    \end{tikzpicture} \end{center}
\end{minipage}
\begin{minipage}[t]{0.35\linewidth}
    \begin{center}
        Earliest Start Time
    \end{center}
\end{minipage}

{~~~}

{~~~}

\begin{minipage}[t]{0.55\linewidth}
    \begin{center} \begin{tikzpicture}[baseline=(current bounding box.center)]
        \definecolor{lightGray}{gray}{0.9}
        \definecolor{darkGray}{gray}{0.7}

        \node[draw=none, fill=lightGray, minimum width=3cm, minimum height=0.5cm] at (0, 0) {};
        \node[draw=none, fill=lightGray, minimum width=3cm, minimum height=0.5cm] at (3.5, 0) {};

        \node[draw=none, fill=darkGray, minimum width=2cm, minimum height=0.5cm] at (1.75, -0.75) {};
    \end{tikzpicture} \end{center}
\end{minipage}
\begin{minipage}[t]{0.35\linewidth}
    \begin{center}
        Shortest Interval
    \end{center}
\end{minipage}

{~~~}

{~~~}

\begin{minipage}[t]{0.55\linewidth}
    \begin{center} \begin{tikzpicture}[baseline=(current bounding box.center)]
        \definecolor{lightGray}{gray}{0.9}
        \definecolor{darkGray}{gray}{0.7}

        \node[draw=none, fill=lightGray, minimum width=1cm, minimum height=0.5cm] at (0, 0) {};
        \node[draw=none, fill=lightGray, minimum width=1cm, minimum height=0.5cm] at (1.5, 0) {};
        \node[draw=none, fill=lightGray, minimum width=1cm, minimum height=0.5cm] at (3, 0) {};
        \node[draw=none, fill=lightGray, minimum width=1cm, minimum height=0.5cm] at (4.5, 0) {};

        \node[draw=none, fill=lightGray, minimum width=1cm, minimum height=0.5cm] at (0.75, -0.75) {};
        \node[draw=none, fill=lightGray, minimum width=1cm, minimum height=0.5cm] at (0.75, -1.5) {};
        \node[draw=none, fill=lightGray, minimum width=1cm, minimum height=0.5cm] at (0.75, -2.25) {};

        \node[draw=none, fill=lightGray, minimum width=1cm, minimum height=0.5cm] at (3.75, -0.75) {};
        \node[draw=none, fill=lightGray, minimum width=1cm, minimum height=0.5cm] at (3.75, -1.5) {};
        \node[draw=none, fill=lightGray, minimum width=1cm, minimum height=0.5cm] at (3.75, -2.25) {};

        \node[draw=none, fill=darkGray, minimum width=1cm, minimum height=0.5cm] at (2.25, -0.75) {};
    \end{tikzpicture} \end{center}
\end{minipage}
\begin{minipage}[t]{0.35\linewidth}
    \begin{center}
        Fewest Conflicts
    \end{center}
\end{minipage}

\subsection{Greedy Algorithm}

We can implement greedy with earliest finish time (EFT)

\begin{listu}
    \item Sort jobs by finish time, say $f_1 \le f_2 \le \dots \le f_n$. \[
        \mathcal{O}(n \log n)
    \]

    \item For each job $j$, we need to check if it's compatible will \itblue{all} previously added jobs. 
    
    \begin{listu}
        \item Natively, this is $\mathcal{O}(n^2)$, as we need $\mathcal{O}(n)$ for each job.
        \item We only need to check if $s_j \ge f_{i^*}$, where $i^*$ is the \itblue{last job added}.
        \begin{listu}
            \item For any jobs $i$ added before $i^*$, we have $f_i \le f_{i^*}$.
            \item By keeping track of $f_{i^*}$, we can check compatibility of job $j$ in $\mathcal{O}(1)$.
        \end{listu}
    \end{listu}

    \item Thus, the total running time is \[ 
        \mathcal{O}(n \log n).
    \]
\end{listu}

\subsection{Proof of Optimality}

\subsubsection{By Contradiction}

\begin{listu}
    \item Suppose for contradiction that greedy is not optimal

    \item Say greedy selects jobs $i_1, i_2, \dots, i_k$ sorted by finish time

    \item Consider an optimal solution $j_1, j_2, \dots, j_m$ sorted by finish time which matches greedy for as many indices as possible. That is, $j_1 = i_1, j_2 = i_2, \dots, j_r = i_r$ for the greatest possible $r$.
    
    \item Both $i_{r+1}$ and $j_{r+1}$ are compatible with $i_1, i_2, \dots, i_r = j_1, j_2, \dots, j_r$.

    \item Consider a new solution $i_1, i_2, \dots, i_r, {\color{red}j_{r+1}}, {\color{lightBlue}j_{r+2}, \dots, j_m}$.
    
    \begin{listu}
        \item We have replaced $j_{r+1}$ by $i_{r+1}$ in our reference optimal solution
        \item This is still \bred{feasible} because $f_{i_{r+1}} \le f_{j_{r+1}} \le s_{j_{t}}$ for $t \ge r+2$.
        \item This is still \bred{optimal} because $m$ jobs are sorted. 
        \item But it matched the greedy solution in $r + 1$ indices. This is a contradiction, as greedy is optimal.
    \end{listu}
\end{listu}

\subsubsection{By Induction}

% TODO: see slides
TODO: SEE SLIDES